\documentclass[10pt]{extarticle}
\usepackage[a4paper,margin=1.5cm]{geometry}
\usepackage[T1]{fontenc}
\usepackage{amsmath,amssymb}
\usepackage{graphicx}
\usepackage[export]{adjustbox}
\usepackage{multicol}
\usepackage{enumitem}
\usepackage{booktabs,array,tabularx}
\setlength{\columnsep}{0.7cm}
\setlength{\parindent}{0pt}
\setlength{\parskip}{3pt}
\pagestyle{empty}
\makeatletter
\newcommand{\qfigdef}[4]{\@namedef{qf@file@#1}{#2}\@namedef{qf@w@#1}{#3}\@namedef{qf@wide@#1}{#4}}
\newcommand{\qf@check}[1]{\@ifundefined{qf@file@#1}{\PackageError{score2}{Unknown figure id #1}{}}{}}
\newcommand{\qf@img}[4][]{\includegraphics[width=#2\dimexpr\@nameuse{qf@w@#4}pt\relax,max width=#3,max height=0.45\textheight,#1]{\@nameuse{qf@file@#4}}}
\newif\ifqf@full
\newif\ifqf@boxed
\newsavebox{\qf@box}
\newcommand{\qf@open}{\global\qf@fulltrue\par\global\qf@boxedtrue
  \setbox\qf@box\vbox\bgroup\hsize\textwidth\linewidth\textwidth}
\newcommand{\qf@emit}[1]{\par\egroup\global\qf@boxedfalse
  \ifdim\dimexpr\ht\qf@box+\dp\qf@box\relax<#1\textheight\noindent\usebox{\qf@box}\par\else\unvbox\qf@box\fi}
\newcommand{\qfullsplit}{\ifqf@boxed\qf@emit{0.85}\fi}
\newcommand{\qf@close}{\ifqf@boxed\qf@emit{0.45}\fi\global\qf@fullfalse}
\newcommand{\qfullbegin}{\end{multicols}\qf@open}
\newcommand{\qfullend}{\qf@close\begin{multicols}{2}}
\newcommand{\qfullstart}{\qf@open}
\newcommand{\qfulldone}{\qf@close}
\newcommand{\qfig}[2][1]{\qf@check{#2}%
  \ifqf@full\qf@blk{#1}{#2}\else
  \ifnum\@nameuse{qf@wide@#2}=1 \expandafter\@firstoftwo\else\expandafter\@secondoftwo\fi
  {\end{multicols}\begin{center}\qf@img{#1}{\textwidth}{#2}\end{center}\begin{multicols}{2}}%
  {\qf@blk{#1}{#2}}\fi}
\newcommand{\qf@blk}[2]{\par\smallskip{\centering\qf@img{#1}{\linewidth}{#2}\par}\smallskip}
\newcommand{\qfigi}[2][1]{\qf@check{#2}\qf@img[valign=c]{#1}{\linewidth}{#2}}
\makeatother
\newcommand{\papersec}[1]{\par\bigskip{\centering\Large\bfseries PAPER-#1\par}\medskip}
\newcommand{\partsec}[2]{\par\medskip\noindent{\large\bfseries PART-#1: #2}\par\smallskip}
\newcommand{\sect}[1]{\par\smallskip\noindent\textbf{#1}\par}
\newcommand{\qn}[1]{\par\smallskip\noindent\textbf{#1.}~}
\newcommand{\paraq}[2]{\par\smallskip\noindent\textbf{Paragraph for Questions #1 and #2}\par}
\newenvironment{qoptions}{\begin{itemize}[leftmargin=1.8em,itemsep=1pt,topsep=2pt,label={}]}{\end{itemize}}
\newcommand{\opt}[2]{\item[(#1)] #2}
% figure registry (generated by score2_tools.py finalize)
\qfigdef{p1_phy_s1_q01_1}{SRG (Enthusiast) - Test 22 (18 Sep 2022)-img/p1_phy_s1_q01_1.png}{129.2}{0}
\qfigdef{p1_phy_s1_q02_1}{SRG (Enthusiast) - Test 22 (18 Sep 2022)-img/p1_phy_s1_q02_1.png}{124.2}{0}
\qfigdef{p1_phy_s1_q03_1}{SRG (Enthusiast) - Test 22 (18 Sep 2022)-img/p1_phy_s1_q03_1.png}{196.5}{0}
\qfigdef{p1_phy_s1_q04_1}{SRG (Enthusiast) - Test 22 (18 Sep 2022)-img/p1_phy_s1_q04_1.png}{91.0}{0}
\qfigdef{p1_phy_s2_q06_A}{SRG (Enthusiast) - Test 22 (18 Sep 2022)-img/p1_phy_s2_q06_A.png}{96.2}{0}
\qfigdef{p1_phy_s2_q06_B}{SRG (Enthusiast) - Test 22 (18 Sep 2022)-img/p1_phy_s2_q06_B.png}{96.2}{0}
\qfigdef{p1_phy_s2_q06_C}{SRG (Enthusiast) - Test 22 (18 Sep 2022)-img/p1_phy_s2_q06_C.png}{89.9}{0}
\qfigdef{p1_phy_s2_q06_D}{SRG (Enthusiast) - Test 22 (18 Sep 2022)-img/p1_phy_s2_q06_D.png}{89.9}{0}
\qfigdef{p1_phy_s2_q08_1}{SRG (Enthusiast) - Test 22 (18 Sep 2022)-img/p1_phy_s2_q08_1.png}{140.3}{0}
\qfigdef{p1_phy_s2_q09_1}{SRG (Enthusiast) - Test 22 (18 Sep 2022)-img/p1_phy_s2_q09_1.png}{124.5}{0}
\qfigdef{p1_phy_s2_q10_1}{SRG (Enthusiast) - Test 22 (18 Sep 2022)-img/p1_phy_s2_q10_1.png}{85.5}{0}
\qfigdef{p1_phy_s2_q13_1}{SRG (Enthusiast) - Test 22 (18 Sep 2022)-img/p1_phy_s2_q13_1.png}{173.9}{0}
\qfigdef{p1_phy_s2_q14_1}{SRG (Enthusiast) - Test 22 (18 Sep 2022)-img/p1_phy_s2_q14_1.png}{117.2}{0}
\qfigdef{p1_phy_s2_q15_1}{SRG (Enthusiast) - Test 22 (18 Sep 2022)-img/p1_phy_s2_q15_1.png}{124.2}{0}
\qfigdef{p1_phy_s3_q02_1}{SRG (Enthusiast) - Test 22 (18 Sep 2022)-img/p1_phy_s3_q02_1.png}{158.2}{0}
\qfigdef{p1_phy_s3_q03_2}{SRG (Enthusiast) - Test 22 (18 Sep 2022)-img/p1_phy_s3_q03_2.png}{169.1}{0}
\qfigdef{p1_phy_s3_q04_1}{SRG (Enthusiast) - Test 22 (18 Sep 2022)-img/p1_phy_s3_q04_1.png}{111.9}{0}
\qfigdef{p1_phy_s3_q05_1}{SRG (Enthusiast) - Test 22 (18 Sep 2022)-img/p1_phy_s3_q05_1.png}{196.5}{0}
\qfigdef{p1_chem_s1_q01_1}{SRG (Enthusiast) - Test 22 (18 Sep 2022)-img/p1_chem_s1_q01_1.png}{167.5}{0}
\qfigdef{p1_chem_s1_q01_A}{SRG (Enthusiast) - Test 22 (18 Sep 2022)-img/p1_chem_s1_q01_A.png}{166.5}{0}
\qfigdef{p1_chem_s1_q01_B}{SRG (Enthusiast) - Test 22 (18 Sep 2022)-img/p1_chem_s1_q01_B.png}{169.5}{0}
\qfigdef{p1_chem_s1_q01_C}{SRG (Enthusiast) - Test 22 (18 Sep 2022)-img/p1_chem_s1_q01_C.png}{152.5}{0}
\qfigdef{p1_chem_s1_q01_D}{SRG (Enthusiast) - Test 22 (18 Sep 2022)-img/p1_chem_s1_q01_D.png}{152.5}{0}
\qfigdef{p1_chem_s1_q02_A}{SRG (Enthusiast) - Test 22 (18 Sep 2022)-img/p1_chem_s1_q02_A.png}{38.2}{0}
\qfigdef{p1_chem_s1_q02_B}{SRG (Enthusiast) - Test 22 (18 Sep 2022)-img/p1_chem_s1_q02_B.png}{28.4}{0}
\qfigdef{p1_chem_s1_q02_C}{SRG (Enthusiast) - Test 22 (18 Sep 2022)-img/p1_chem_s1_q02_C.png}{28.4}{0}
\qfigdef{p1_chem_s1_q02_D}{SRG (Enthusiast) - Test 22 (18 Sep 2022)-img/p1_chem_s1_q02_D.png}{28.5}{0}
\qfigdef{p1_chem_s1_q05_A}{SRG (Enthusiast) - Test 22 (18 Sep 2022)-img/p1_chem_s1_q05_A.png}{55.6}{0}
\qfigdef{p1_chem_s1_q05_B}{SRG (Enthusiast) - Test 22 (18 Sep 2022)-img/p1_chem_s1_q05_B.png}{74.2}{0}
\qfigdef{p1_chem_s1_q05_C}{SRG (Enthusiast) - Test 22 (18 Sep 2022)-img/p1_chem_s1_q05_C.png}{74.4}{0}
\qfigdef{p1_chem_s1_q05_D}{SRG (Enthusiast) - Test 22 (18 Sep 2022)-img/p1_chem_s1_q05_D.png}{65.6}{0}
\qfigdef{p1_chem_s2_q06_1}{SRG (Enthusiast) - Test 22 (18 Sep 2022)-img/p1_chem_s2_q06_1.png}{53.3}{0}
\qfigdef{p1_chem_s2_q06_2}{SRG (Enthusiast) - Test 22 (18 Sep 2022)-img/p1_chem_s2_q06_2.png}{66.4}{0}
\qfigdef{p1_chem_s2_q06_3}{SRG (Enthusiast) - Test 22 (18 Sep 2022)-img/p1_chem_s2_q06_3.png}{78.7}{0}
\qfigdef{p1_chem_s2_q06_4}{SRG (Enthusiast) - Test 22 (18 Sep 2022)-img/p1_chem_s2_q06_4.png}{69.8}{0}
\qfigdef{p1_chem_s2_q07_A}{SRG (Enthusiast) - Test 22 (18 Sep 2022)-img/p1_chem_s2_q07_A.png}{93.2}{0}
\qfigdef{p1_chem_s2_q07_B}{SRG (Enthusiast) - Test 22 (18 Sep 2022)-img/p1_chem_s2_q07_B.png}{78.1}{0}
\qfigdef{p1_chem_s2_q07_C}{SRG (Enthusiast) - Test 22 (18 Sep 2022)-img/p1_chem_s2_q07_C.png}{76.5}{0}
\qfigdef{p1_chem_s2_q07_D}{SRG (Enthusiast) - Test 22 (18 Sep 2022)-img/p1_chem_s2_q07_D.png}{77.0}{0}
\qfigdef{p1_chem_s2_q07_5}{SRG (Enthusiast) - Test 22 (18 Sep 2022)-img/p1_chem_s2_q07_5.png}{79.2}{0}
\qfigdef{p1_chem_s2_q08_A}{SRG (Enthusiast) - Test 22 (18 Sep 2022)-img/p1_chem_s2_q08_A.png}{56.7}{0}
\qfigdef{p1_chem_s2_q08_B}{SRG (Enthusiast) - Test 22 (18 Sep 2022)-img/p1_chem_s2_q08_B.png}{63.1}{0}
\qfigdef{p1_chem_s2_q08_D}{SRG (Enthusiast) - Test 22 (18 Sep 2022)-img/p1_chem_s2_q08_D.png}{106.5}{0}
\qfigdef{p1_chem_s2_q09_A}{SRG (Enthusiast) - Test 22 (18 Sep 2022)-img/p1_chem_s2_q09_A.png}{166.9}{0}
\qfigdef{p1_chem_s2_q09_B}{SRG (Enthusiast) - Test 22 (18 Sep 2022)-img/p1_chem_s2_q09_B.png}{200.5}{0}
\qfigdef{p1_chem_s2_q09_C}{SRG (Enthusiast) - Test 22 (18 Sep 2022)-img/p1_chem_s2_q09_C.png}{216.0}{0}
\qfigdef{p1_chem_s2_q09_D}{SRG (Enthusiast) - Test 22 (18 Sep 2022)-img/p1_chem_s2_q09_D.png}{155.5}{0}
\qfigdef{p1_chem_s2_q11_1}{SRG (Enthusiast) - Test 22 (18 Sep 2022)-img/p1_chem_s2_q11_1.png}{185.5}{0}
\qfigdef{p1_chem_s2_q11_A}{SRG (Enthusiast) - Test 22 (18 Sep 2022)-img/p1_chem_s2_q11_A.png}{77.0}{0}
\qfigdef{p1_chem_s2_q11_B}{SRG (Enthusiast) - Test 22 (18 Sep 2022)-img/p1_chem_s2_q11_B.png}{28.8}{0}
\qfigdef{p1_chem_s2_q11_C}{SRG (Enthusiast) - Test 22 (18 Sep 2022)-img/p1_chem_s2_q11_C.png}{60.6}{0}
\qfigdef{p1_chem_s2_q11_D}{SRG (Enthusiast) - Test 22 (18 Sep 2022)-img/p1_chem_s2_q11_D.png}{60.6}{0}
\qfigdef{p1_chem_s2_q12_1}{SRG (Enthusiast) - Test 22 (18 Sep 2022)-img/p1_chem_s2_q12_1.png}{122.6}{0}
\qfigdef{p1_chem_s2_q12_A}{SRG (Enthusiast) - Test 22 (18 Sep 2022)-img/p1_chem_s2_q12_A.png}{86.2}{0}
\qfigdef{p1_chem_s2_q12_B}{SRG (Enthusiast) - Test 22 (18 Sep 2022)-img/p1_chem_s2_q12_B.png}{67.0}{0}
\qfigdef{p1_chem_s2_q12_C}{SRG (Enthusiast) - Test 22 (18 Sep 2022)-img/p1_chem_s2_q12_C.png}{94.0}{0}
\qfigdef{p1_chem_s2_q12_D}{SRG (Enthusiast) - Test 22 (18 Sep 2022)-img/p1_chem_s2_q12_D.png}{47.4}{0}
\qfigdef{p1_chem_s2_q13_1}{SRG (Enthusiast) - Test 22 (18 Sep 2022)-img/p1_chem_s2_q13_1.png}{56.3}{0}
\qfigdef{p1_chem_s2_q13_A}{SRG (Enthusiast) - Test 22 (18 Sep 2022)-img/p1_chem_s2_q13_A.png}{45.8}{0}
\qfigdef{p1_chem_s2_q13_B}{SRG (Enthusiast) - Test 22 (18 Sep 2022)-img/p1_chem_s2_q13_B.png}{69.8}{0}
\qfigdef{p1_chem_s2_q13_C}{SRG (Enthusiast) - Test 22 (18 Sep 2022)-img/p1_chem_s2_q13_C.png}{45.8}{0}
\qfigdef{p1_chem_s2_q13_D}{SRG (Enthusiast) - Test 22 (18 Sep 2022)-img/p1_chem_s2_q13_D.png}{61.5}{0}
\qfigdef{p1_chem_s2_q14_1}{SRG (Enthusiast) - Test 22 (18 Sep 2022)-img/p1_chem_s2_q14_1.png}{85.1}{0}
\qfigdef{p1_chem_s2_q14_A}{SRG (Enthusiast) - Test 22 (18 Sep 2022)-img/p1_chem_s2_q14_A.png}{99.9}{0}
\qfigdef{p1_chem_s2_q14_B}{SRG (Enthusiast) - Test 22 (18 Sep 2022)-img/p1_chem_s2_q14_B.png}{34.9}{0}
\qfigdef{p1_chem_s2_q14_C}{SRG (Enthusiast) - Test 22 (18 Sep 2022)-img/p1_chem_s2_q14_C.png}{34.9}{0}
\qfigdef{p1_chem_s2_q14_D}{SRG (Enthusiast) - Test 22 (18 Sep 2022)-img/p1_chem_s2_q14_D.png}{49.5}{0}
\qfigdef{p1_chem_s2_q15_A}{SRG (Enthusiast) - Test 22 (18 Sep 2022)-img/p1_chem_s2_q15_A.png}{134.2}{0}
\qfigdef{p1_chem_s2_q15_B}{SRG (Enthusiast) - Test 22 (18 Sep 2022)-img/p1_chem_s2_q15_B.png}{129.6}{0}
\qfigdef{p1_chem_s2_q15_C}{SRG (Enthusiast) - Test 22 (18 Sep 2022)-img/p1_chem_s2_q15_C.png}{116.7}{0}
\qfigdef{p1_chem_s2_q15_D}{SRG (Enthusiast) - Test 22 (18 Sep 2022)-img/p1_chem_s2_q15_D.png}{118.9}{0}
\qfigdef{p1_chem_s3_q01_1}{SRG (Enthusiast) - Test 22 (18 Sep 2022)-img/p1_chem_s3_q01_1.png}{236.5}{0}
\qfigdef{p1_chem_s3_q03_1}{SRG (Enthusiast) - Test 22 (18 Sep 2022)-img/p1_chem_s3_q03_1.png}{75.5}{0}
\qfigdef{p1_chem_s3_q03_3}{SRG (Enthusiast) - Test 22 (18 Sep 2022)-img/p1_chem_s3_q03_3.png}{91.6}{0}
\qfigdef{p1_chem_s3_q03_2}{SRG (Enthusiast) - Test 22 (18 Sep 2022)-img/p1_chem_s3_q03_2.png}{103.5}{0}
\qfigdef{p1_chem_s3_q03_5}{SRG (Enthusiast) - Test 22 (18 Sep 2022)-img/p1_chem_s3_q03_5.png}{93.2}{0}
\qfigdef{p1_chem_s3_q03_4}{SRG (Enthusiast) - Test 22 (18 Sep 2022)-img/p1_chem_s3_q03_4.png}{61.3}{0}
\qfigdef{p1_chem_s3_q03_7}{SRG (Enthusiast) - Test 22 (18 Sep 2022)-img/p1_chem_s3_q03_7.png}{82.5}{0}
\qfigdef{p1_chem_s3_q03_6}{SRG (Enthusiast) - Test 22 (18 Sep 2022)-img/p1_chem_s3_q03_6.png}{76.8}{0}
\qfigdef{p1_chem_s3_q03_8}{SRG (Enthusiast) - Test 22 (18 Sep 2022)-img/p1_chem_s3_q03_8.png}{67.6}{0}
\qfigdef{p1_chem_s3_q03_9}{SRG (Enthusiast) - Test 22 (18 Sep 2022)-img/p1_chem_s3_q03_9.png}{80.1}{0}
\qfigdef{p1_chem_s3_q03_10}{SRG (Enthusiast) - Test 22 (18 Sep 2022)-img/p1_chem_s3_q03_10.png}{81.4}{0}
\qfigdef{p1_chem_s3_q03_11}{SRG (Enthusiast) - Test 22 (18 Sep 2022)-img/p1_chem_s3_q03_11.png}{91.0}{0}
\qfigdef{p1_chem_s3_q05_1}{SRG (Enthusiast) - Test 22 (18 Sep 2022)-img/p1_chem_s3_q05_1.png}{90.7}{0}

\begin{document}
\begin{center}{\Large\bfseries SRG (Enthusiast) --- Test 22 (18 Sep 2022)}\end{center}
\begin{multicols}{2}
\partsec{1}{PHYSICS}
\sect{Section-I(i)}
\qn{1} The circuit was in the shown state for a long time. Now if the switch S is closed then the net charge that flows through the switch S, will be
\qfig{p1_phy_s1_q01_1}
\begin{qoptions}
\opt{A}{$\dfrac{400}{3}\,\mu$C}
\opt{B}{$100\,\mu$C}
\opt{C}{$\dfrac{100}{3}\,\mu$C}
\opt{D}{$50\,\mu$C}
\end{qoptions}
\qn{2} A parallel plate capacitor of capacitance C (without dielectrics) is filled by dielectric slabs as shown in figure. Then the new capacitance of the capacitor is:
\qfig{p1_phy_s1_q02_1}
\begin{qoptions}
\opt{A}{3.9 C}
\opt{B}{4 C}
\opt{C}{2.4 C}
\opt{D}{3 C}
\end{qoptions}
\qn{3} A uniform right-angled triangular lamina is suspended from the ceiling with the help of three light strings. The plane of the lamina is horizontal and the strings are vertical as shown in the figure. In which of the following options, tensile forces $T_A$, $T_B$ and $T_C$ developed in the strings connected to the vertices A, B and C respectively of the lamina are arranged in correct order of their magnitudes?
\qfig{p1_phy_s1_q03_1}
\begin{qoptions}
\opt{A}{$T_A>T_B>T_C$}
\opt{B}{$T_B>T_A>T_C$}
\opt{C}{$T_C>T_A>T_B$}
\opt{D}{$T_A=T_B=T_C$}
\end{qoptions}
\qn{4} A uniform steel rod leans against a frictionless wall in static equilibrium. Frictional force between the lower end and the floor is less than its limiting value by a finite amount. The rod is supplied some amount of heat so that it expands. Assume that the coefficient of friction does not change on heating. What can you conclude regarding contact forces from the wall and from the floor?
\qfig{p1_phy_s1_q04_1}
\begin{qoptions}
\opt{A}{From the wall it decreases and from the floor it increases.}
\opt{B}{From the wall it decreases and from the floor it decreases.}
\opt{C}{From the wall as well as from the floor it increases.}
\opt{D}{From the wall as well as from the floor it decreases.}
\end{qoptions}
\qn{5} A parallel plate capacitor consists of square plates of edge length $a$ separated by a distance $d$ that is much smaller that the edge length. It is charged to a potential $V$ and made to move with a constant velocity $v$ directed along one of its edges. How much magnetic field exists inside the capacitor?
\begin{qoptions}
\opt{A}{$\dfrac{\mu_0\varepsilon_0\,vV}{d}$}
\opt{B}{$\dfrac{\mu_0\varepsilon_0\,vV}{a}$}
\opt{C}{$\dfrac{\mu_0\varepsilon_0\,dvV}{a^2}$}
\opt{D}{None of these}
\end{qoptions}
\sect{Section-I(ii)}
\qn{6} In which of the following system(s) linear momentum cannot be conserved during collision along horizontal line (parallel to base)\,?
\begin{qoptions}
\opt{A}{\qfigi{p1_phy_s2_q06_A}}
\opt{B}{\qfigi{p1_phy_s2_q06_B}}
\opt{C}{\qfigi{p1_phy_s2_q06_C}}
\opt{D}{\qfigi{p1_phy_s2_q06_D}}
\end{qoptions}
\qn{7} A particle moving with kinetic energy $=3$ J makes an elastic head-on collision with a stationary particle which has twice its mass. During the impact:
\begin{qoptions}
\opt{A}{The minimum kinetic energy of the system is 1 J.}
\opt{B}{The maximum elastic potential energy of the system is 2 J.}
\opt{C}{Momentum and total energy are conserved at every instant.}
\opt{D}{All of the above are correct}
\end{qoptions}
\qn{8} Two smooth balls A and B, each of mass $m$ and radius $R$, have their centres at $(0,0,R)$ and at $(5R,-R,R)$ respectively, in a coordinate system as shown. Ball A, moving along positive x-axis, collides with ball B. Just before the collision, speed of ball A is 4 m/s and ball B is stationary. The collision between the balls is elastic.
\qfig{p1_phy_s2_q08_1}
\begin{qoptions}
\opt{A}{Velocity of the ball A just after the collision is $\left(\hat i+\sqrt3\,\hat j\right)$ m/s}
\opt{B}{Velocity of B just after collision is zero.}
\opt{C}{Impulse of the force exerted by A on B during the collision, is equal to $\left(3m\hat i-\sqrt3 m\hat j\right)$ kg--m/s}
\opt{D}{Impulse of the force exerted by A on B is $\left(4m\hat i\right)$ kg--m/s}
\end{qoptions}
\qn{9} A rigid square loop of side length $\ell$ carrying an electric current is held motionless on a frictionless horizontal tabletop. A uniform magnetic field pointing upwards is switched on everywhere to the left of the dashed line as shown in the figure and then the loop is released. Considering different lengths of the side segments x and y ($x<y$), which of the following conclusions can you make?
\qfig{p1_phy_s2_q09_1}
\begin{qoptions}
\opt{A}{If $x<\dfrac{\ell}{2}$ and $y=\dfrac{\ell}{2}$, the loop starts rotating anticlockwise}
\opt{B}{If $x<\dfrac{\ell}{2}$ and $y<\dfrac{\ell}{2}$, the loop starts rotating anticlockwise}
\opt{C}{If $x<\dfrac{\ell}{2}$ and $y>\dfrac{\ell}{2}$, the loop starts rotating clockwise}
\opt{D}{If $x<\dfrac{\ell}{2}$ and $y>\dfrac{\ell}{2}$, more information is required to decide which way the loop starts rotating}
\end{qoptions}
\qn{10} A positively charged particle P crosses a region of uniform electric field of intensity E. Velocity $u$ of the particle at the entry on the left side of the region is perpendicular to the electric field as shown in the first figure. Speed of the particle at exit on the right side is $v_1$. Now the experiment is repeated with a uniform magnetic field of induction B superimposed on the electric field as shown in the second figure, keeping entry velocity of the particle unchanged. In this experiment, speed of the particle at exit on the right side is $v_2$. What can you conclude regarding the speeds $v_1$ and $v_2$?
\qfig{p1_phy_s2_q10_1}
\begin{qoptions}
\opt{A}{If $u<\dfrac{E}{B}$, then $v_2<v_1$}
\opt{B}{If $u=\dfrac{E}{B}$, then $v_2<v_1$}
\opt{C}{If $u=\dfrac{E}{B}$, then $v_2=v_1$}
\opt{D}{If $u>\dfrac{E}{B}$, then $v_2<v_1$}
\end{qoptions}
\qn{11} A human body has a surface area of approximately $1\,\mathrm{m}^2$. The normal body temperature is 10 K above the surrounding room temperature $T_0$. Take the room temperature to be $T_0=300$ K. For $T_0=300$ K, the value of $\sigma T_0^4=460\ \mathrm{Wm^{-2}}$ (where $\sigma$ is the Stefan-Boltzmann constant). Which of the following options is/are correct?
\begin{qoptions}
\opt{A}{The amount of energy radiated by the body in 1 second is close to 60 Joules.}
\opt{B}{If the surrounding temperature reduces by a small amount $\Delta T_0\ll T_0$, then to maintain the same body temperature the same (living) human being needs to radiate $\Delta Q=4\sigma T_0^3\Delta T_0$ more energy per unit time}
\opt{C}{Reducing the exposed surface area of the body (e.g. by curling up) allows humans to maintain the same body temperature while reducing the energy lost by radiation.}
\opt{D}{If the body temperature rises significantly then the peak in the spectrum of electromagnetic radiation emitted by the body would shift to longer wavelengths.}
\end{qoptions}
\qn{12} You are a consultant for a cookware manufacturer who wishes to make a pan that will have two features:\\
(1) Absorb thermal energy from a flame as quickly as possible, and\\
(2) Have an inner surface that remains as hot as possible when cooking.\\
Consider only radiation. You should recommend a pan with the
\begin{qoptions}
\opt{A}{Outer and inner surfaces black.}
\opt{B}{Outer and inner surfaces shiny.}
\opt{C}{Outer surface shiny and inner surface black.}
\opt{D}{Outer surface black and inner surface shiny.}
\end{qoptions}
\qn{13} Two identical small beads A and B threaded on an almost inertia-less horizontal rod of length $2\ell$ are connected by a light inextensible cord of length $1.2\ell$. Initially the beads are adjusted equidistant from the centre of the rod and the rod is rotating about its central vertical axis at an angular velocity $1/3$ rad/s. At some point of time, the cord is cut and the beads begin to slide on the rod without friction. Neglect effects of gravity. Which of the following statements are correct?
\qfig{p1_phy_s2_q13_1}
\begin{qoptions}
\opt{A}{Path of a bead in an inertial frame is a spiral.}
\opt{B}{Path of a bead in an inertial frame is a straight line.}
\opt{C}{To leave the rod, the beads take 4s after the cord is cut.}
\opt{D}{To leave the rod, the beads take more than 4s after the cord is cut.}
\end{qoptions}
\qn{14} An iron block and a wooden block are positioned in a vessel containing water as shown in the figure. The iron block (1) hangs from a massless string with a rigid support from the top while the wooden block (2) floats being tied to the bottom through a massless string. If now the vessel starts acceleration towards right.
\qfig{p1_phy_s2_q14_1}
\begin{qoptions}
\opt{A}{Iron block gets deflected towards right.}
\opt{B}{Wooden block gets deflected towards right.}
\opt{C}{Iron block gets deflected towards left.}
\opt{D}{Wooden block gets deflected towards left.}
\end{qoptions}
\qn{15} Water jet is projected at an angle to the horizontal. At the point of projection, the area of the jet is $S_1$ and at the highest point, the area of the jet is $S_2$. The initial velocity of projection is $u$.
\qfig{p1_phy_s2_q15_1}
\begin{qoptions}
\opt{A}{The angle of projection is $\cos^{-1}\!\left(\dfrac{S_1}{S_2}\right)$}
\opt{B}{The range on the level ground is $\dfrac{2u^2}{g}\dfrac{S_1}{S_2}\sqrt{1-\dfrac{S_1^2}{S_2^2}}$}
\opt{C}{The maximum height reached from the ground is $\dfrac{2u^2}{g}\left(1-\dfrac{S_1^2}{S_2^2}\right)$}
\opt{D}{The rate of volume flow is $S_2u$}
\end{qoptions}
\sect{Section-II}
\qn{1} Three small balls A, B and C each of mass $m$ are connected by three thin light rods each of length $\ell$ to make an equilateral triangular structure. The structure is sliding on a frictionless horizontal floor with all the three balls on the floor. At an instant, the ball A is observed moving with a velocity $v$ along the rod AB and the ball B moving along the rod BC. Find moduli of tensile forces developed in rod AB [in N] [Given: $m=1$ kg, $v=2$ m/s, $\ell=2$ m]
\qn{2} Shows the top view of a cylindrical can mounted on a turntable. The can is filled with water. At a depth $h$ below the water surface are two horizontal tubes of length $\ell$ and cross-sectional area $a$, with right-angle bends at their ends. As the water jets emerge from the tubes, there is a torque $\tau$ exerted on the system given by the expression $\tau=x\rho gh(r+\ell)a$, where $\rho$ is the density of the water. Write `$x$'.
\qfig{p1_phy_s3_q02_1}
\qn{3} Consider the arrangement shown, where distance between the plates of a parallel plate capacitor is $d$. When a potential difference $V$ is applied between the plates, electrons are emitted from the negative plate with negligible initial velocity. These electrons are attracted towards the positive plate creating a small electric current, which is registered in a micro-ammeter. What minimum magnetic field must be applied perpendicular to the electric field in the capacitor so that the micro-ammeter registers no current? Denote charge and mass of an electron by symbols $e$ and $m$ respectively and neglect effects of gravity. If minimum magnetic field $\dfrac{\mu\,1}{2\,d}\sqrt{\dfrac{nmv}{e}}$. Find `$n$'
\qfig{p1_phy_s3_q03_2}
\qn{4} Three rods AO, BO and CO having equal lengths and equal area of cross section are connected as shown in the figure. The temperature of ends A, B and C are $30^\circ$C, $60^\circ$C and $90^\circ$ respectively. The thermal conductivities of rods AO, BO and CO are in the ratio of $1:2:3$. What is the temperature of junction O? [in $^\circ$C]
\qfig{p1_phy_s3_q04_1}
\qn{5} Two identical plates of metal are welded end to end as shown in figure (A), 20 cal of heat flows through it in 4 minutes when the ends were maintained at different temperatures. If the plates are welded as shown in figure (B) and the ends were maintain at same temperature as in first case, the same amount of heat will flow through the plates in how many minutes\,?
\qfig{p1_phy_s3_q05_1}
\partsec{2}{CHEMISTRY}
\sect{Section-I(i)}
\qn{1} Identify P and Q respectively
\qfig{p1_chem_s1_q01_1}
\begin{qoptions}
\opt{A}{\qfigi{p1_chem_s1_q01_A}}
\opt{B}{\qfigi{p1_chem_s1_q01_B}}
\opt{C}{\qfigi{p1_chem_s1_q01_C}}
\opt{D}{\qfigi{p1_chem_s1_q01_D}}
\end{qoptions}
\qn{2} The most stable carbocation is
\begin{qoptions}
\opt{A}{\qfigi{p1_chem_s1_q02_A}}
\opt{B}{\qfigi{p1_chem_s1_q02_B}}
\opt{C}{\qfigi{p1_chem_s1_q02_C}}
\opt{D}{\qfigi{p1_chem_s1_q02_D}}
\end{qoptions}
\qn{3} Calculate the electronegativity of X (Bond energy of AX bond is 78 KCal/mole). Given that bond energy of A$_2$ and X$_2$ are 36 and 49 KCal/mol respectively. Given that Electronegativity of A is 3
\begin{qoptions}
\opt{A}{2.752}
\opt{B}{1.752}
\opt{C}{3.752}
\opt{D}{0.752}
\end{qoptions}
\qn{4} How many dibromoderivatives of 2, 2, 5, 5 - tetramethyl hexane are there (excluding stereoisomers)
\begin{qoptions}
\opt{A}{6}
\opt{B}{7}
\opt{C}{8}
\opt{D}{9}
\end{qoptions}
\qn{5} Which of the following is most reactive toward HCN\,?
\begin{qoptions}
\opt{A}{\qfigi{p1_chem_s1_q05_A}}
\opt{B}{\qfigi{p1_chem_s1_q05_B}}
\opt{C}{\qfigi{p1_chem_s1_q05_C}}
\opt{D}{\qfigi{p1_chem_s1_q05_D}}
\end{qoptions}
\sect{Section-I(ii)}
\qn{6} Compound P upon reductive ozonolysis produce Q and R respectively. P upon Hydrogenation with H$_2$/Pd gives 2, 4 - dimethyl hexane. Q undergoes haloform reaction but not Tollen's test. R undergoes crossed aldol reaction with HCHO in presence of dil.NaOH gives only S. S reacts with NaBH$_4$/CH$_3$OH gives U. U reacts with HCHO/H$^+$ gives V. Then find correct options\,?
\begin{qoptions}
\opt{A}{V is \qfigi{p1_chem_s2_q06_1}}
\opt{B}{P is \qfigi{p1_chem_s2_q06_2}}
\opt{C}{S is \qfigi{p1_chem_s2_q06_3}}
\opt{D}{U is \qfigi{p1_chem_s2_q06_4}}
\end{qoptions}
\qn{7} Among the following compounds, optically inactive is/are
\begin{qoptions}
\opt{A}{\qfigi{p1_chem_s2_q07_A}}
\opt{B}{\qfigi{p1_chem_s2_q07_B}}
\opt{C}{\qfigi{p1_chem_s2_q07_C}}
\opt{D}{\qfigi{p1_chem_s2_q07_D}}
\end{qoptions}
\qn{8} Which of the following more acidic than \qfigi{p1_chem_s2_q07_5}
\begin{qoptions}
\opt{A}{\qfigi{p1_chem_s2_q08_A}}
\opt{B}{\qfigi{p1_chem_s2_q08_B}}
\opt{C}{HCOOH}
\opt{D}{\qfigi{p1_chem_s2_q08_D}}
\end{qoptions}
\columnbreak\qn{9} Which of the following has pair of identical molecules
\begin{qoptions}
\opt{A}{\qfigi{p1_chem_s2_q09_A}}
\opt{B}{\qfigi{p1_chem_s2_q09_B}}
\opt{C}{\qfigi{p1_chem_s2_q09_C}}
\opt{D}{\qfigi{p1_chem_s2_q09_D}}
\end{qoptions}
\qn{10} Which of the following is/are correctly matched
\begin{qoptions}
\opt{A}{Na$^+>$ Mg$^{+2}>$ Al$^{+3}>$ O$^{-2}$ (radius)}
\opt{B}{Cr $>$ Mn $>$ V $>$ Sc (I.E.$_2$)}
\opt{C}{Cs$_2$O $>$ BaO $>$ MgO $>$ Al$_2$O$_3$ [Basic strength]}
\opt{D}{O $>$ S $>$ Se $>$ Te [Electron gain enthalpy]}
\end{qoptions}
\qn{11} \qfig{p1_chem_s2_q11_1}
In the above sequence the final products of the reaction is/are\,:
\begin{qoptions}
\opt{A}{\qfigi{p1_chem_s2_q11_A}}
\opt{B}{\qfigi{p1_chem_s2_q11_B}}
\opt{C}{\qfigi{p1_chem_s2_q11_C}}
\opt{D}{\qfigi{p1_chem_s2_q11_D}}
\end{qoptions}
\columnbreak\qn{12} \qfig{p1_chem_s2_q12_1}
Possible options for A\,?
\begin{qoptions}
\opt{A}{\qfigi{p1_chem_s2_q12_A}}
\opt{B}{\qfigi{p1_chem_s2_q12_B}}
\opt{C}{\qfigi{p1_chem_s2_q12_C}}
\opt{D}{\qfigi{p1_chem_s2_q12_D}}
\end{qoptions}
\qn{13} Which of the following is/are more stable than \qfigi{p1_chem_s2_q13_1}
\begin{qoptions}
\opt{A}{\qfigi{p1_chem_s2_q13_A}}
\opt{B}{\qfigi{p1_chem_s2_q13_B}}
\opt{C}{\qfigi{p1_chem_s2_q13_C}}
\opt{D}{\qfigi{p1_chem_s2_q13_D}}
\end{qoptions}
\qn{14} Which of the following is/are less basic than \qfigi{p1_chem_s2_q14_1}
\begin{qoptions}
\opt{A}{\qfigi{p1_chem_s2_q14_A}}
\opt{B}{\qfigi{p1_chem_s2_q14_B}}
\opt{C}{\qfigi{p1_chem_s2_q14_C}}
\opt{D}{\qfigi{p1_chem_s2_q14_D}}
\end{qoptions}
\qn{15} In which of the following reaction product gives intramolecular aldol condensation reaction\,?
\begin{qoptions}
\opt{A}{\qfigi{p1_chem_s2_q15_A}}
\opt{B}{\qfigi{p1_chem_s2_q15_B}}
\opt{C}{\qfigi{p1_chem_s2_q15_C}}
\opt{D}{\qfigi{p1_chem_s2_q15_D}}
\end{qoptions}
\sect{Section-II}
\qn{1} Consider the reaction sequence from P to T.
\qfig{p1_chem_s3_q01_1}
The yield of reaction P to Q 100 \%\\
The yield of reaction Q to R 50\%\\
The yield of reaction R to S 25\%\\
The yield of reaction S to T 100\%\\
What is the amount in grams of T obtained from 216 grams of P.\\
{[}Molar mass of C = 12, H = 1, O = 16{]}
\qn{2} Ionization potential of element X is 1200 kJ/mole and electron gain enthalpy of element X is $-350$ kJ/mole. Then find electronegativity of X in Mulliken Scale [1 eV/atom = 96.35 kJ/mole] \{Answer rounded of to nearest integer\}
\columnbreak\qn{3} How many of the following more acidic than \qfigi{p1_chem_s3_q03_1}\\
(1) CH$_3$COOH \hfill (2) HCOOH
\par\smallskip
\noindent\begin{tabular}{@{}p{0.49\linewidth}@{}p{0.49\linewidth}@{}}
\qfigi{p1_chem_s3_q03_3} & \qfigi{p1_chem_s3_q03_2} \\[2pt]
\qfigi{p1_chem_s3_q03_5} & \qfigi{p1_chem_s3_q03_4} \\[2pt]
\qfigi{p1_chem_s3_q03_7} & \qfigi{p1_chem_s3_q03_6} \\[2pt]
\qfigi{p1_chem_s3_q03_8} & \qfigi{p1_chem_s3_q03_9} \\[2pt]
\qfigi{p1_chem_s3_q03_10} & \qfigi{p1_chem_s3_q03_11}
\end{tabular}
\qn{4} Compound (X) having M.F ($\mathrm{C_5H_{10}O}$) [only ketone and aldehyde] reacts with\\
(i) C$_2$H$_5$MgBr/(ii) H$_3$O$^+$ gives alcohol. Find no. of X which can give alcohol with at least one chiral carbon.
\qn{5} Number of dibromoderivatives of biphenyl \qfigi{p1_chem_s3_q05_1} (only structural isomers)
\partsec{3}{MATHEMATICS}
\sect{Section-I(i)}
\qn{1} The value of the integral $I=\displaystyle\int_2^3\sqrt{(x-2)(3-x)}\,dx$ is equal to
\begin{qoptions}
\opt{A}{$\dfrac{\pi}{4}$}
\opt{B}{$\dfrac{\pi}{2}$}
\opt{C}{$\dfrac{\pi}{8}$}
\opt{D}{$2\pi$}
\end{qoptions}
\qn{2} Let $x>0$ be a fixed real number. Then the integral $\displaystyle\int_0^\infty e^{-t}\,|x-t|\,dt$ is equal to
\begin{qoptions}
\opt{A}{$x+2e^{-x}-1$}
\opt{B}{$x-2e^{-x}+1$}
\opt{C}{$x+2e^{-x}+1$}
\opt{D}{$-x-2e^{-x}+1$}
\end{qoptions}
\qn{3} If $f(x)=xf^2(x)+x^2\ \forall x\in\mathbb R$, then $\displaystyle\int\dfrac{f^2(x)-2xf(x)+2x+1}{\bigl(1-2xf(x)\bigr)f(x)\bigl(f(x)(1-x)+2x+1\bigr)}\,dx$ equals
\begin{qoptions}
\opt{A}{$\dfrac{-1}{x-f(x)}+c$}
\opt{B}{$\dfrac{-1}{x+f(x)}+c$}
\opt{C}{$\dfrac{-1}{x^2-f(x)}+c$}
\opt{D}{$\dfrac{-1}{x^2+f(x)}+c$}
\end{qoptions}
\qn{4} $\displaystyle\int\left(x^{24}+x^{16}+x^8\right)\left(2x^{16}+3x^8+6\right)^{1/8}dx=\dfrac{\left(2x^{24}+3x^{16}+6x^8\right)^{\beta/\gamma}}{\alpha}+c$ (where $\beta,\gamma$ are co-prime natural numbers and $c$ be constant of integration), then the value of $\alpha+\beta+\gamma$ is
\begin{qoptions}
\opt{A}{61}
\opt{B}{51}
\opt{C}{71}
\opt{D}{91}
\end{qoptions}
\qn{5} $\displaystyle\int\dfrac{(2x+1)}{\left(x^2+4x+1\right)^{3/2}}\,dx$
\begin{qoptions}
\opt{A}{$\dfrac{x^3}{\left(x^2+4x+1\right)^{1/2}}+C$}
\opt{B}{$\dfrac{x}{\left(x^2+4x+1\right)^{1/2}}+C$}
\opt{C}{$\dfrac{x^2}{\left(x^2+4x+1\right)^{1/2}}+C$}
\opt{D}{$\dfrac{1}{\left(x^2+4x+1\right)^{1/2}}+C$}
\end{qoptions}
\sect{Section-I(ii)}
\qn{6} Consider an equation, $y\left(3^{|x|}-1\right)+2|x|\left(2^{y}-1\right)=0$ where $y=\log_2\left(3x^2-1\right)-\log_2\sqrt{x^2+1}$. Identify which of the following statement(s) is (are) correct\,?
\begin{qoptions}
\opt{A}{Number of real solutions of the equation is 2}
\opt{B}{Number of real solutions of the equation is 5}
\opt{C}{Sum of squares of all the solutions is $\dfrac79$}
\opt{D}{Sum of squares of all the solutions is $\dfrac{14}{9}$}
\end{qoptions}
\qn{7} $f:[0,1]\to\mathbb R$ be a function such that $f^2(x)=1$ and $f(x)$ is discontinuous at exactly two point and $\displaystyle\int_0^1 f(x)\,dx=0$. Let $I=\displaystyle\int_0^1 xf(x)\,dx$, choose the incorrect option(s)
\begin{qoptions}
\opt{A}{Maximum value of $I$ is $\dfrac14$}
\opt{B}{Minimum value of $I$ is $-\dfrac34$}
\opt{C}{When $I$ is maximum then number of such function $f(x)$ is 8}
\opt{D}{When $I$ is minimum then number of such function $f(x)$ is 4}
\end{qoptions}
\qn{8} If $\displaystyle\int_{-\alpha}^{\alpha}\Bigl\{e^x+(\cos x)\bigl(\ln\bigl(x+\sqrt{1+x^2}\bigr)\bigr)\Bigr\}dx>\dfrac32$, then $\alpha$ can be
\begin{qoptions}
\opt{A}{1}
\opt{B}{2}
\opt{C}{3}
\opt{D}{4}
\end{qoptions}
\qn{9} If $\displaystyle\int_0^1\left\{\left(1-x^{2018}\right)^{\frac1{2020}}-\left(1-x^{2020}\right)^{\frac1{2018}}\right\}dx=J$, then
\begin{qoptions}
\opt{A}{$J<\dfrac{2018}{2020}$}
\opt{B}{$J>\dfrac{2018}{2020}$}
\opt{C}{$J>\dfrac{2020}{2018}$}
\opt{D}{$J<\dfrac{2020}{2018}$}
\end{qoptions}
\qn{10} If $\displaystyle\int\left(2x^6+15x^4+2x^2+3\right)\cos(2x)\,dx=f(x)\sin(2x)+g(x)\cos(2x)+K$, where $f(x)$ and $g(x)$ are polynomial of $x$ such that $f(0)=1$, $g(-1)=-4$ and $K$ is constant of integration, then
\begin{qoptions}
\opt{A}{$4f(1)\le 3g(1)$}
\opt{B}{$\displaystyle\int_{-\pi/3}^{\pi/3}f(x)\,dx\ne0$}
\opt{C}{$g''(0)<10$}
\opt{D}{$\dfrac{g(2)}{f(2)}>1$}
\end{qoptions}
\qn{11} Which of the following function(s) is/are even\,?
\begin{qoptions}
\opt{A}{$f(x)=\displaystyle\int_0^x\ln\left(t+\sqrt{1+t^2}\right)dt$}
\opt{B}{$g(x)=\displaystyle\int_0^x\dfrac{\left(2^t+1\right)t}{2^t-1}\,dt$}
\opt{C}{$h(x)=\displaystyle\int_0^x\left(\sqrt{1+t+t^2}-\sqrt{1-t+t^2}\right)dt$}
\opt{D}{$\ell(x)=\displaystyle\int_0^x\ln\left(\dfrac{1-t}{1+t}\right)dt$}
\end{qoptions}
\qn{12} If $\displaystyle\int\dfrac{3x\sin^2x\cos x-3\sin^3x}{x^4}\,dx=f(x)+c$, $\displaystyle\lim_{x\to0}f(x)=1$, then which of the following is/are true.
\begin{qoptions}
\opt{A}{$\displaystyle\lim_{x\to0}\dfrac{\int_0^x tf(t)\,dt-2x^2}{1-\cos x}=-3$}
\opt{B}{$\displaystyle\lim_{x\to0}\dfrac{\left(f(x)\right)^{\frac13}-x^2}{x^2}=-\dfrac16$}
\opt{C}{If $h(x)=x\cdot\sqrt[3]{f(x)}$, then $\displaystyle\int_0^\pi h^4(x)\,dx=\dfrac{3\pi}{8}$}
\opt{D}{If $h(x)=x\cdot\sqrt[3]{f(x)}$, then {\footnotesize$\displaystyle\int_0^{\pi/2}e^{h(x)}\left(\cos^2x-\sin x\right)dx=-1$}}
\end{qoptions}
\qn{13} If $f:(0,1)\to\mathbb R$ is defined by $f(x)=\displaystyle\sum_{n=0}^{\infty}\dfrac{x^{2^n}}{1-x^{2^{n+1}}}$, then --
\begin{qoptions}
\opt{A}{$f^{-1}(x)=\dfrac{x}{x+1}$}
\opt{B}{$\displaystyle\int_0^{1/2}f(x)\,dx=\dfrac{2\ln2-1}{2}$}
\opt{C}{$\displaystyle\int_0^1f(x)\,dx=\ln2-1$}
\opt{D}{$\displaystyle\int_0^1f(x)\,dx=\dfrac12+\dfrac13+\dfrac14+\dots\infty$}
\end{qoptions}
\qn{14} Let $I=\displaystyle\int_0^1\sqrt{\dfrac{1+\sqrt x}{1-\sqrt x}}\,dx$ and $J=\displaystyle\int_0^1\sqrt{\dfrac{1-\sqrt x}{1+\sqrt x}}\,dx$ then which statement(s) are correct
\begin{qoptions}
\opt{A}{$I+J=2$}
\opt{B}{$I-J=\pi$}
\opt{C}{$J=\dfrac{2+\pi}{2}$}
\opt{D}{$J=\dfrac{4-\pi}{2}$}
\end{qoptions}
\qn{15} $\displaystyle\int\dfrac{x^4+1}{x^6+1}\,dx$ is
\begin{qoptions}
\opt{A}{$\tan^{-1}x+\dfrac13\tan^{-1}\left(x^3\right)+c$}
\opt{B}{$\dfrac13\tan^{-1}\left(\dfrac{3x-3x^5}{1-3x^2-3x^4+x^6}\right)+c$}
\opt{C}{$\dfrac13\tan^{-1}\left(\dfrac{2x-x^4}{1-2x^5-3x^2}\right)+c$}
\opt{D}{$\dfrac12\tan^{-1}\left(\dfrac{x-x^3-1}{1+x^2-3x^5}\right)+c$}
\end{qoptions}
\sect{Section-II}
\qn{1} Let $f(1)+g(1)=9e$; $f(x)=-x^2g'(x)$ and $g(x)=-x^2f'(x)$. If $\displaystyle\int_1^4\left(\dfrac{f(x)+g(x)}{x^2}\right)dx=k\left(e-e^{\frac14}\right)$, then $\left|\dfrac k2\right|$ is
\qn{2} If $f\bigl(x+f(y)\bigr)+f\bigl(x-f(y)\bigr)=f(x)\ \forall\, x,y\in\mathbb R$ and $(0,2)$ lies on $y=f(x)$, then $\displaystyle\int_0^{12}f(x)\,dx=\lambda\int_{-24}^{24}f(x-36)\,dx$ where $\lambda$ equals
\qn{3} $\displaystyle\int_{\pi/6}^{\pi/3}\dfrac{2022-2023\csc^2x}{\sin^{2022}x}\,dx$ equals $-\dfrac{2^p\left(3^q-1\right)}{3^r}$, then $p+q+r$ is equal to
\qn{4} If $\displaystyle\int\dfrac{\left(\sqrt x\right)^5}{\left(\sqrt x\right)^7+x^6}\,dx=a\cdot\ln\left(\dfrac{x^{k/2}}{x^{k/2}+1}\right)+C$, then the value of $k$ is (given $a>0$)
\qn{5} If $A=\{a_1,a_2,\dots,a_n\}$ be set of all possible natural factors of 1000 and $P=\displaystyle\int_{-2}^{2}x\ln\left(a_1^x+a_2^x+\dots+a_n^x\right)dx$, then value of $\dfrac{2P}{n}$ equals $\ln k$, find $k$.
\end{multicols}
\end{document}
